% Pendahuluan buku. Bagian ini memuat ciri karya ilmiah utuh yang dituntut
% sebuah buku referensi, yaitu latar belakang, rumusan masalah, metode
% penyusunan, serta kebaruan, lalu menjelaskan cara memakai buku ini sebagai
% buku ajar.
\chapter{Pendahuluan}
\label{ch:pendahuluan}

\section*{Latar Belakang}
\addcontentsline{toc}{section}{Latar Belakang}

Arsitektur perangkat lunak adalah susunan bagian besar sebuah sistem beserta
aturan hubungan di antaranya. Pembedanya dari rancangan rinci adalah biaya
perubahan. Keputusan yang menyentuh hampir seluruh \textit{file} termasuk
keputusan arsitektur, sehingga keputusan itu jarang dapat ditarik kembali dengan
murah. Karena itu keputusan arsitektur sebaiknya diambil dengan dasar yang dapat
diperiksa.

Buku acuan utama bab-bab di sini, yaitu karya Richards dan Ford, menyediakan
dasar itu dalam bentuk penilaian~\cite{richards2020fundamentals}. Setiap gaya
arsitektur diberi peringkat bintang untuk sejumlah karakteristik, misalnya
\textit{scalability}, kinerja, dan kemudahan pengujian. Peringkat seperti itu
berguna sebagai peta awal. Namun peringkat tidak menyebut seberapa besar selisih
antara dua gaya, dan tidak dapat diulang di komputer pembaca.

Buku ini menempuh jalan yang berbeda. Setiap klaim tentang sebuah gaya
dirumuskan menjadi klaim yang dapat dibuktikan salah, atau \textit{falsifiable
claim}, lalu diperiksa dengan program Python kecil yang disertakan. Klaim
\textit{ports and adapters}, misalnya, menyatakan bahwa \textit{core logic}
tidak mengenal teknologi \textit{storage} apa pun. Satu baris impor sudah cukup
untuk menggugurkan klaim itu, sehingga pemeriksaannya cukup menghitung baris
impor.

Klaim yang menyangkut waktu menuntut kehati-hatian lebih, sebab selisihnya
sering kecil. Pengukurannya diulang sedikitnya seratus putaran, lalu dilaporkan
sebagai median, persentil ke-95, dan nilai terbesar. Selisih dua pilihan
kemudian dibandingkan dengan \textit{range} pengukurannya sendiri. Pada Bab~2,
misalnya, jalan pintas \textit{open layer} hanya menghemat 0,0046 milidetik pada
median, lebih kecil daripada \textit{range}-nya. Pencacahan per putaran tetap
menunjukkan jalan pintas itu lebih cepat pada 97 dari 100 putaran.

Contoh pada buku ini diambil dari keadaan yang dikenal pembaca di Indonesia.
Peralihan layanan transfer antarbank ke BI-FAST dibahas pada Bab~1. Kurs dasar
pada Bab~2 meniru kurs referensi Jakarta Interbank Spot Dollar Rate (JISDOR)
dari Bank Indonesia. Bab lain memakai urusan kampus, misalnya pembayaran uang
kuliah, pengisian rencana studi, dan dompet mahasiswa.

\section*{Arti Kata Terukur}
\addcontentsline{toc}{section}{Arti Kata Terukur}

Judul buku ini memakai kata terukur\index{terukur} untuk cara menilai yang
diuraikan pada Latar Belakang. Sebuah klaim arsitektur disebut terukur bila
memenuhi tiga syarat. Pertama, klaimnya menyebut besaran beserta satuannya,
misalnya jumlah penyebutan, milidetik, atau \textit{message} per
\textit{lookup}. Kedua, angkanya keluar dari program di \berkas{sources/} yang
dijalankan menurut prosedur tertulis, yaitu langkah kerja pada Latihan. Ketiga,
angka itu dapat diperoleh ulang oleh pembaca di komputernya sendiri.

Analoginya adalah konsumsi bensin sepeda motor. Iklan yang hanya menyebut irit
tidak dapat diperiksa, sebab setiap orang menafsirkan kata itu berbeda. Uji jalan
yang menyebut jarak tempuh dalam kilometer per liter pada rute tertentu dapat
diulang siapa pun. Angkanya bergeser bila jalannya lain, tetapi dua motor yang
diuji pada rute yang sama tetap dapat dibandingkan. Angka pada buku ini pun
begitu. Milidetiknya bergeser di komputer lain, sedangkan perbandingan dua versi
yang dijalankan di komputer yang sama dapat diperiksa ulang.

Tabel~\ref{tab:klaim-terukur} memperlihatkan bedanya pada lima bab. Kolom kedua
memuat klaim yang lazim diucapkan tanpa angka. Kolom ketiga memuat bentuk
terukurnya beserta hasil pengukuran pada bab tersebut.

\begin{table}[htbp]
  \centering
  \caption{Klaim yang lazim diucapkan tanpa angka dan bentuk terukurnya pada
    lima bab}
  \label{tab:klaim-terukur}
  \begin{tabularx}{\textwidth}{|c|>{\raggedright\arraybackslash}p{3.3cm}%
    |>{\raggedright\arraybackslash}X|}
    \hline
    \textbf{Bab} & \textbf{Klaim tanpa angka} &
      \textbf{Bentuk terukur dan hasilnya} \\
    \hline
    1 & Server sanggup melayani banyak klien & Median naik dari 0,31 menjadi
      1,36 ms pada enam belas klien, tetapi nilai terbesarnya melampaui
      1.000 ms sejak delapan klien \\
    \hline
    2 & \textit{Open layer} membuat \textit{request} lebih cepat & Jalan pintas
      menghemat 0,0046 ms pada median, lebih kecil daripada \textit{range}
      0,0376 ms, tetapi lebih cepat pada 97 dari 100 putaran \\
    \hline
    4 & \textit{Core logic} terlepas dari basis data & Nol teknologi
      \textit{storage} disebut \textit{core logic}, dan berpindah dari memori
      ke SQLite mengubah nol \textit{file} \\
    \hline
    9 & Event membuat \textit{request} lebih cepat & Latensi ujung ke ujung
      justru 6,2 kali pada topologi broker dan 9,3 kali pada topologi mediator,
      dengan susunan sinkron menang pada 200 dari 200 putaran \\
    \hline
    14 & \textit{Serverless} lebih murah & Pada 720 \textit{request} per jam,
      tagihannya 20,88 melawan 1.800,00 gigabita detik, dengan titik impas
      62.069 \textit{request} per jam \\
    \hline
  \end{tabularx}
\end{table}

Bentuk terukur juga dapat menggugurkan klaimnya sendiri. Baris Bab~9
memperlihatkannya, sebab anggapan bahwa event mempercepat \textit{request} gugur
pada seluruh 200 putaran. Klaim yang tidak dapat digugurkan oleh angka apa pun
belum terukur.

Kata terukur tidak berarti angkanya berlaku di semua komputer. Angka milidetik
pada buku ini terikat pada perangkat di Tabel~\ref{tab:lingkungan-ukur}, dan
yang dapat dipakai ulang adalah pola perbandingannya. Kata itu juga tidak
berarti keputusan arsitektur cukup diambil dari angka. Pada Bab~2, jalan pintas
yang terbukti lebih cepat tetap ditolak untuk kurs jual, sebab tuntutan bisnisnya
menyatakan hanya staf yang boleh membacanya.

Dalam pemakaian sehari-hari, kata terukur juga berarti diperhitungkan, misalnya
pada frasa kebijakan yang terukur. Arti kedua itu ikut dimaksud judul buku ini.
Keputusan arsitektur menjadi terukur ketika harga setiap pilihan sudah diketahui
sebelum pilihan itu diambil, bukan sesudah sistemnya berjalan.

\section*{Rumusan Masalah}
\addcontentsline{toc}{section}{Rumusan Masalah}

Latar belakang tersebut menghasilkan tiga pertanyaan yang dijawab buku ini.

\begin{enumerate}
  \item Bagaimana klaim kualitas sebuah gaya arsitektur dirumuskan menjadi
        \textit{falsifiable claim} yang dapat diperiksa dengan program kecil?
  \item Berapa harga yang dibayar setiap gaya untuk kelebihan yang
        dijanjikannya, bila harga itu diukur dengan pengulangan dan sebaran?
  \item Kapan selisih yang terukur cukup kuat untuk menentukan pilihan gaya,
        dan kapan selisih itu tenggelam di dalam ragam pengukurannya sendiri?
\end{enumerate}

Kebaruan pertanyaan tersebut terletak pada bentuk jawabannya. Peringkat pada
buku acuan menjawab gaya mana yang lebih baik untuk sebuah karakteristik.
Pertanyaan di atas menuntut jawaban berupa angka, beserta cara mengulang angka
itu di komputer lain. Ketiganya dijawab kembali pada bagian Penutup.

\section*{Tujuan dan Manfaat}
\addcontentsline{toc}{section}{Tujuan dan Manfaat}

Tujuan buku ini adalah menyediakan rujukan berbahasa Indonesia tentang empat
belas gaya arsitektur. Harga setiap gaya dinyatakan dalam angka yang dapat
diukur ulang, dan setiap angka disertai program yang menghasilkannya.

Manfaatnya berbeda menurut pembacanya. Pengembang dan arsitek memperoleh dasar
untuk menimbang \textit{trade-off} sebelum keputusan yang mahal diubah diambil.
Peneliti memperoleh prosedur pengukuran dan kode yang dapat diperluas, misalnya
untuk sistem yang memanggil model bahasa. Pengajar memperoleh bahan yang
langkah kerja, tabel contoh, dan lembar isiannya sudah tersedia.

\section*{Ruang Lingkup dan Batasan}
\addcontentsline{toc}{section}{Ruang Lingkup dan Batasan}

Bab~1 sampai Bab~14 masing-masing membahas satu gaya atau satu kelompok gaya,
mulai dari client-server sampai \textit{serverless}. Kelompok gaya itu misalnya
empat varian Model-View-* pada Bab~3, serta Command Query Responsibility
Segregation (CQRS) bersama \textit{event sourcing} pada Bab~10. Nama lengkap gaya
pada Bab~12 adalah \textit{orchestration-driven service-oriented architecture}
(SOA). Seluruh gaya itu kemudian diterapkan pada satu sistem yang sudah berjalan,
yaitu Exim, asisten tanya jawab ekspor impor pada Bab~15.

Batasannya ada tiga. Pertama, seluruh program ditulis dalam Python, dan sebagian
besar hanya memakai pustaka standar. Kedua, seluruh pengukuran berjalan pada
satu laptop di alamat lokal, sehingga waktu tempuh jaringan antarkota tidak
termuat. Ketiga, skala contohnya kecil, misalnya \textit{processing unit} pada
Bab~8 yang diukur pada satu, empat, dan tujuh unit. Angkanya karena itu
menunjukkan pola perbandingan, bukan kapasitas sistem produksi.

\section*{Metode Penyusunan}
\addcontentsline{toc}{section}{Metode Penyusunan}

Setiap bab disusun melalui enam langkah yang sama.

\begin{enumerate}
  \item Klaim tentang gaya yang dibahas diambil dari rujukan primer, yaitu buku
        acuan, spesifikasi resmi, atau dokumentasi lembaga yang membakukannya.
  \item Klaim itu dirumuskan menjadi \textit{falsifiable claim} yang menyebut
        angka, misalnya nol penyebutan atau selisih dalam milidetik.
  \item Program Python kecil ditulis dalam dua versi atau lebih yang
        mengerjakan pekerjaan sama, misalnya \textit{closed layer} dan
        \textit{open layer}. Selisih antarversi karena itu berasal dari
        susunannya.
  \item Program dijalankan berulang. Selisih yang kecil diukur sedikitnya
        seratus putaran, lalu dilaporkan median, persentil ke-95, nilai
        terbesar, dan \textit{win count} setiap pilihan.
  \item Keluaran dikutip sebagaimana tercetak, termasuk angka yang janggal.
        Sebab kejanggalannya dijelaskan pada paragraf sesudahnya.
  \item Kode yang tampil di naskah disalin dari \textit{file} yang dijalankan.
        Potongannya boleh lebih pendek, tetapi tidak boleh berbeda isinya.
\end{enumerate}

Seluruh pengukuran dijalankan pada satu laptop dengan perangkat pada
Tabel~\ref{tab:lingkungan-ukur}. Aplikasi contoh dan alat ukurnya berjalan di
komputer yang sama. Versi Python yang dipakai juga tercetak pada keluaran
Listing~\ref{lst:venv-ch01} di Bab~1.

\begin{table}[htbp]
  \centering
  \caption{Perangkat yang dipakai menjalankan pengukuran di buku ini}
  \label{tab:lingkungan-ukur}
  \begin{tabularx}{\textwidth}{|l|X|}
    \hline
    \textbf{Perangkat} & \textbf{Keterangan} \\
    \hline
    Prosesor & Intel Core i9-11900H, delapan inti, enam belas \textit{thread} \\
    \hline
    Memori & 30 GiB, sebagaimana terbaca oleh sistem operasi \\
    \hline
    Sistem operasi & Ubuntu 24.04 dengan kernel Linux 6.8 \\
    \hline
    Penerjemah & Python 3.14.10 di dalam venv \\
    \hline
    Pustaka tambahan & Tkinter pada Bab~3, Pillow pada Bab~6, serta Flask dan
      requests pada Bab~11 sampai Bab~13 \\
    \hline
  \end{tabularx}
\end{table}

Angka waktu pada buku ini karena itu terikat pada komputer tersebut. Nilai
milidetiknya akan berbeda di komputer lain. Yang dapat dipakai ulang adalah
pola perbandingannya, misalnya pilihan mana yang lebih cepat dan pada berapa
dari seratus putaran.

\section*{Kebaruan dan Kontribusi}
\addcontentsline{toc}{section}{Kebaruan dan Kontribusi}

Kebaruan buku ini terletak pada cara setiap gaya dinilai. Penilaiannya tidak
berhenti pada uraian dan peringkat, tetapi diteruskan sampai angka yang
dihasilkan program yang disertakan. Aturan pengukurannya sama pada keempat belas
bab, sehingga pola perbandingan antarbab dapat dibaca berdampingan.

Kontribusi yang menyertainya berupa program Python per bab di dalam
\berkas{sources/} dan 42 prosedur pengukuran yang disertakan sebagai Latihan.
Setiap prosedur punya tabel contoh yang diisi dari pengukuran nyata, sehingga
hasil pembaca dapat dibandingkan langsung. Hasil yang berlawanan dengan anggapan
umum pun dicetak apa adanya. \textit{Event-driven architecture}, misalnya, tidak
mempercepat satu \textit{request}, sebab latensinya 6,2 sampai 9,3 kali latensi
susunan sinkron.

\section*{Sistematika Penulisan}
\addcontentsline{toc}{section}{Sistematika Penulisan}

Tabel~\ref{tab:peta-buku} mendaftar gaya yang dibahas setiap bab beserta tiga
hal yang diukur pada Latihan-nya. Bab~1 sampai Bab~14 mengikuti urutan yang
sama, yaitu Tujuan Pembelajaran, Pendahuluan, seksi materi, \textit{Summary},
lalu Latihan. Bab~15 tidak punya seksi Latihan, sebab seluruh isinya adalah
tugas proyek.

\begin{table}[htbp]
  \centering
  \caption{Gaya yang dibahas setiap bab beserta hal yang diukur}
  \label{tab:peta-buku}
  \begin{tabularx}{\textwidth}{|c|>{\raggedright\arraybackslash}p{3.6cm}%
    |>{\raggedright\arraybackslash}X|}
    \hline
    \textbf{Bab} & \textbf{Judul bab} & \textbf{Hal yang diukur pada Latihan} \\
    \hline
    1 & Pengenalan Arsitektur Perangkat Lunak dan Client-Server &
      Tanggung jawab klien dan server, biaya perjalanan, jumlah klien \\
    \hline
    2 & Layered Architecture & Jalan satu \textit{request}, \textit{layer
      violation}, biaya satu \textit{hop} \\
    \hline
    3 & Model-View-* & Arah \textit{dependency}, pengujian tanpa antarmuka,
      luas dampak perubahan \\
    \hline
    4 & Ports and Adapters & \textit{Core independence}, \textit{adapter swap
      cost}, \textit{testing in isolation} \\
    \hline
    5 & Microkernel dan Plugin & \textit{Core independence}, \textit{error
      isolation}, \textit{activation cost} \\
    \hline
    6 & Pipeline dan Pipe-and-Filter & \textit{Filter independence},
      \textit{stage reordering}, \textit{first output latency} \\
    \hline
    7 & Peer-to-Peer & \textit{Lookup cost}, \textit{failure resilience},
      \textit{message scalability} \\
    \hline
    8 & Space-Based Architecture & \textit{Throughput scaling}, \textit{stale
      read} beserta \textit{data collision}, \textit{data loss} \\
    \hline
    9 & Event-Driven Architecture & \textit{Latency cost}, \textit{failure
      containment}, \textit{change cost} \\
    \hline
    10 & CQRS dan Event Sourcing & \textit{Write cost}, \textit{replay cost},
      \textit{temporal query} \\
    \hline
    11 & Service-Based Architecture & Tambahan \textit{hop}, satu layanan yang
      mati, perubahan skema basis data bersama \\
    \hline
    12 & Orchestration-Driven SOA & Tambahan \textit{hop}, penyisipan satu
      langkah, perubahan layanan yang dipakai banyak alur kerja \\
    \hline
    13 & Microservices & Latensi dan jumlah \textit{hop}, harga penggabungan,
      rilis ulang satu layanan \\
    \hline
    14 & Serverless Architecture & \textit{Cold start cost}, \textit{idle cost},
      \textit{concurrency limit} \\
    \hline
    15 & Capstone Project & Satu dimensi kualitas pada satu fitur Exim,
      diperbaiki lalu dibuktikan dengan angka \\
    \hline
  \end{tabularx}
\end{table}

Sesudah Bab~15, bagian Penutup menjawab kembali ketiga pertanyaan pada Rumusan
Masalah. Glosarium memuat istilah yang dibakukan di seluruh bab. Indeks menunjuk
halaman tempat setiap istilah kunci pertama kali muncul pada tiap bab.

\section*{Penggunaan sebagai Buku Ajar}
\addcontentsline{toc}{section}{Penggunaan sebagai Buku Ajar}

Buku ini adalah buku referensi. Susunannya juga memungkinkan buku ini digunakan
sebagai buku ajar mata kuliah Arsitektur Perangkat Lunak (IF231303) pada Program
Studi Teknologi Informasi, Universitas Pradita. Bagian yang melayani penggunaan
itu sudah tersedia, dan Tabel~\ref{tab:bagian-ajar} mendaftar letaknya.

\begin{table}[htbp]
  \centering
  \caption{Bagian buku yang melayani penggunaan sebagai buku ajar}
  \label{tab:bagian-ajar}
  \begin{tabularx}{\textwidth}{|>{\raggedright\arraybackslash}X%
    |>{\raggedright\arraybackslash}X|l|}
    \hline
    \textbf{Bagian} & \textbf{Letak} & \textbf{Jumlah} \\
    \hline
    Tujuan Pembelajaran & Awal Bab~1 sampai Bab~15 & Tiga butir per bab \\
    \hline
    \textit{Summary} & Akhir seksi materi setiap bab & 15 \\
    \hline
    Latihan & Akhir Bab~1 sampai Bab~14 & 42 \\
    \hline
    Tabel contoh yang sudah terisi & Di dalam setiap Latihan & 42 \\
    \hline
    Lembar isian & Di dalam setiap Latihan & 42 \\
    \hline
    Proyek akhir beserta pilihan judulnya & Bab~15 & 20 judul \\
    \hline
    Slide & \berkas{slides/sessionNN/} & 15 set \\
    \hline
    Kode dan skrip pengukuran & \berkas{sources/chapterNN/} & 15 direktori \\
    \hline
  \end{tabularx}
\end{table}

Pembaca yang memakai buku ini sebagai rujukan cukup membaca seksi materi dan
\textit{Summary}. Latihan tetap berguna bagi pembaca tersebut, sebab langkah
kerja dan tabel contohnya adalah prosedur untuk mengulang angka pada bab yang
sama. Lembar isian, pertanyaan analisis, dan Bab~15 ditujukan bagi penggunaan di
kelas.

% Nomor tabel terakhir disimpan agar tabel Penutup melanjutkannya.
\setcounter{tabelluarbab}{\value{table}}
